% Slide 10 — Entradas do AIPE
\begin{frame}{Entradas do AIPE}
  \begin{columns}[T,onlytextwidth]
    \column{0.33\textwidth}
      \begin{block}{Dados preparados}
        \small Séries loja-produto, ciclos comerciais (17/ano,
        $\approx$21 dias), filtragem e limpeza.
      \end{block}
    \column{0.33\textwidth}
      \begin{block}{Características operacionais}
        \small 17 variáveis: intermitência, distribuição, tendência,
        sazonalidade, volume.
      \end{block}
    \column{0.33\textwidth}
      \begin{block}{Sinais de previsão}
        \small MASE, sMAPE, RMSSE de 4 modelos (Croston, ARIMA, XGBoost,
        LSTM).
      \end{block}
  \end{columns}
  \vfill
  \begin{alertblock}{Atenção}
    \small Caracterização e previsão são \highlight{entradas}, não a
    responsabilidade final do AIPE. A previsão é \highlight{sinal auxiliar},
    não critério final de escolha de política.
  \end{alertblock}
  \note{
    É importante posicionar corretamente essas três entradas. Os dados
    preparados padronizam séries loja-produto a partir do histórico
    transacional. As características operacionais resumem o comportamento
    da demanda em 17 variáveis numéricas. E a previsão de demanda entra como
    sinal auxiliar -- ela ajuda a caracterizar a previsibilidade da série,
    mas não decide isoladamente qual política usar. Essa distinção evita um
    erro comum: tratar previsão e caracterização como se fossem o produto
    final do AIPE, quando na verdade elas alimentam o núcleo de seleção.
  }
\end{frame}
